% TOP_PROBLEM: {"id": 1242, "title": "Lehmer's Totient Problem", "category_id": 1, "category": {"id": 1, "name": "number_theory", "display_name": "Number Theory", "description": "Properties of integers, prime numbers, Diophantine equations.", "slug": "number-theory", "order_index": 1, "created_at": "2026-07-31T15:26:25.670Z"}, "status": "open"}
% FIRST_POSTED: null
% ENTRY_KIND: statement_only
% Imported from: batch05.tex; UnsolvedMath ID 1242
\documentclass[11pt]{article}
\usepackage[margin=25mm]{geometry}
\usepackage{fontspec}
\setmainfont{Latin Modern Roman}
\usepackage{amsmath,amssymb,mathtools}
\usepackage{unicode-math}
\setmathfont{Latin Modern Math}
\usepackage{xurl,hyperref}
\hypersetup{colorlinks=true,urlcolor=blue}
\setcounter{secnumdepth}{0}
\setlength{\parindent}{0pt}
\setlength{\parskip}{5pt}
\setlength{\emergencystretch}{3em}
\newcommand{\sref}[2]{\hyperlink{source:#1:#2}{[S#2]}}
\newcommand{\cataloguescope}[1]{\paragraph{Recorded target.}#1}
\begin{document}
% UnsolvedMath ID 1242; source catalogue code NT-035
\setcounter{section}{1241}
\section{Lehmer's Totient Problem}\label{1242}
{\small\sffamily UnsolvedMath ID 1242 · Number Theory}
\subsection{Definitions and mathematical statement}

\paragraph{Shared notation.}$\mathbb N=\{1,2,\ldots\}$, and $\mathbb Z,\mathbb Q,\mathbb R,\mathbb C$ denote the integers, rationals, reals and complex numbers. Any other conventions are specified locally.


For an integer $n\geq1$, Euler\textquotesingle s totient is $\varphi(n)=\#\{a\in\{1,\ldots,n\}:\gcd(a,n)=1\}$. An integer is composite if it is greater than $1$ and not prime. In the standard Lehmer problem the domain is $n>1$; with the convention $\varphi(1)=1$, omitting this domain would introduce the unrelated trivial exception $n=1$.
\paragraph{Statement recorded in the supplied source.}
For every integer $n>1$, does
\[
\varphi(n)\mid n-1\quad\Longrightarrow\quad n\text{ is prime}
\]
hold? Equivalently, is there no composite integer $n$ satisfying $\varphi(n)\mid n-1$? \sref{1242}{2}

\subsection{Short English statement}
Can a composite integer have its totient divide one less than itself?
\subsection{Sources}
\begin{description}
\item[\hypertarget{source:1242:1}{S1}] ULAM.AI and the UnsolvedMath Contributors, UnsolvedMath: A Curated Collection of Open Mathematics Problems, record 1242 (NT-035). CC BY 4.0. \url{https://huggingface.co/datasets/ulamai/UnsolvedMath}.
\item[\hypertarget{source:1242:2}{S2}] Manuel Norman, \emph{A group-theoretical approach to Lehmer\textquotesingle s totient problem}, arXiv:2106.11781v1 (2021), abstract and Section 1. \url{https://arxiv.org/pdf/2106.11781}.
\end{description}
\label{end:1242}
\end{document}
